\documentclass[11pt,a4paper]{article}

\usepackage[utf8]{inputenc}
\usepackage[T1]{fontenc}
\usepackage{amsmath,amssymb,amsthm}
% booktabs not available; use \hline instead
\usepackage{geometry}
\usepackage{hyperref}
% enumitem not available; using default list spacing
\usepackage{array}

\geometry{margin=2.5cm}

\hypersetup{
  colorlinks=true,
  linkcolor=blue,
  citecolor=blue,
  urlcolor=blue,
}

\newcommand{\bk}{\mathbf{k}}
\newcommand{\by}{\mathbf{y}}
\newcommand{\bx}{\mathbf{x}}
\newcommand{\br}{\mathbf{r}}
\newcommand{\bff}{\mathbf{f}}
\newcommand{\bigO}{\mathcal{O}}
\newcommand{\dt}{\Delta t}

\title{Numerical Integration Methods\\[4pt]
       \large Reference Documentation for \texttt{rk\_integrators.h}}
\author{Pendulum Simulation Project}
\date{\today}

\begin{document}
\maketitle
\tableofcontents
\newpage

%=============================================================================
\section{Introduction}
%=============================================================================

This document provides a mathematical reference for the numerical integration
methods implemented in \texttt{modules/rk\_integrators.h}. The methods fall into
four families:
\begin{enumerate}
  \item \textbf{Explicit Runge--Kutta methods} (RK3, RK4, RK5, RK23, RKF45)
  \item \textbf{Symplectic integrators} (Semi-Implicit Euler, Leapfrog, Ruth4)
  \item \textbf{Second-order ODE specialists} (Velocity Verlet, Runge--Kutta--Nystr\"{o}m, Numerov)
  \item \textbf{Exponential integrators} (DEN3)
\end{enumerate}

Throughout this document the ODE system is either the first-order form
\begin{equation}\label{eq:first-order}
  \frac{d}{dt}\by = \bff(t,\by),
  \qquad
  \by = \begin{pmatrix}\theta \\ \omega\end{pmatrix},
\end{equation}
or the direct second-order form
\begin{equation}\label{eq:second-order}
  \ddot\theta = a(t,\theta).
\end{equation}
The time step is~$h$ (code variable \texttt{dt}), and subscript~$n$ denotes the
value at time~$t_n$.

%=============================================================================
\section{Classical Runge--Kutta (4th Order) --- \texttt{rk4\_step}}
\label{sec:rk4}
%=============================================================================

The classical four-stage, fourth-order explicit Runge--Kutta method. It is the
most widely used single-step ODE integrator and provides a good balance between
accuracy and computational cost.

\subsection{Formulation}
\begin{align}
  \bk_1 &= \bff(t_n,\;\by_n) \notag\\
  \bk_2 &= \bff\!\left(t_n + \tfrac{h}{2},\;\by_n + \tfrac{h}{2}\,\bk_1\right) \notag\\
  \bk_3 &= \bff\!\left(t_n + \tfrac{h}{2},\;\by_n + \tfrac{h}{2}\,\bk_2\right) \notag\\
  \bk_4 &= \bff\!\left(t_n + h,\;\by_n + h\,\bk_3\right) \notag\\[6pt]
  \by_{n+1} &= \by_n + \frac{h}{6}\bigl(\bk_1 + 2\,\bk_2 + 2\,\bk_3 + \bk_4\bigr)
  \label{eq:rk4}
\end{align}

\subsection{Butcher Tableau}
\[
\renewcommand{\arraystretch}{1.3}
\begin{array}{c|cccc}
  0   &      &      &      &     \\
  1/2 & 1/2  &      &      &     \\
  1/2 & 0    & 1/2  &      &     \\
  1   & 0    & 0    & 1    &     \\
  \hline
      & 1/6  & 1/3  & 1/3  & 1/6 \\
\end{array}
\]

\subsection{Properties}
\begin{center}
\begin{tabular}{ll}
  \hline
  Order          & 4 \\
  Stages         & 4 \\
  Implicit       & No \\
  Symplectic     & No \\
  Local error    & $\bigO(h^5)$ \\
  \hline
\end{tabular}
\end{center}

%=============================================================================
\section{Standard Runge--Kutta (3rd Order) --- \texttt{rk3\_step}}
\label{sec:rk3}
%=============================================================================

A three-stage, third-order explicit Runge--Kutta method. Lower cost per step
than RK4 but also lower accuracy; useful for convergence-order studies.

\subsection{Formulation}
\begin{align}
  \bk_1 &= \bff(t_n,\;\by_n) \notag\\
  \bk_2 &= \bff\!\left(t_n + \tfrac{h}{2},\;\by_n + \tfrac{h}{2}\,\bk_1\right) \notag\\
  \bk_3 &= \bff\!\left(t_n + h,\;\by_n - h\,\bk_1 + 2h\,\bk_2\right) \notag\\[6pt]
  \by_{n+1} &= \by_n + \frac{h}{6}\bigl(\bk_1 + 4\,\bk_2 + \bk_3\bigr)
  \label{eq:rk3}
\end{align}

\subsection{Butcher Tableau}
\[
\renewcommand{\arraystretch}{1.3}
\begin{array}{c|ccc}
  0   &      &     &     \\
  1/2 & 1/2  &     &     \\
  1   & -1   & 2   &     \\
  \hline
      & 1/6  & 2/3 & 1/6 \\
\end{array}
\]

\subsection{Properties}
\begin{center}
\begin{tabular}{ll}
  \hline
  Order          & 3 \\
  Stages         & 3 \\
  Implicit       & No \\
  Symplectic     & No \\
  Local error    & $\bigO(h^4)$ \\
  \hline
\end{tabular}
\end{center}

%=============================================================================
\section{Butcher's Runge--Kutta (5th Order) --- \texttt{rk5\_step}}
\label{sec:rk5}
%=============================================================================

A six-stage, fifth-order explicit Runge--Kutta method due to Butcher~(1964).
Higher accuracy per step than RK4 at the cost of two additional derivative
evaluations.

\subsection{Formulation}
\begin{align}
  \bk_1 &= \bff(t_n,\;\by_n) \notag\\
  \bk_2 &= \bff\!\left(t_n + \tfrac{h}{4},\;\by_n + \tfrac{h}{4}\,\bk_1\right) \notag\\
  \bk_3 &= \bff\!\left(t_n + \tfrac{h}{4},\;\by_n + \tfrac{h}{8}\,\bk_1
            + \tfrac{h}{8}\,\bk_2\right) \notag\\
  \bk_4 &= \bff\!\left(t_n + \tfrac{h}{2},\;\by_n - \tfrac{h}{2}\,\bk_2
            + h\,\bk_3\right) \notag\\
  \bk_5 &= \bff\!\left(t_n + \tfrac{3h}{4},\;\by_n + \tfrac{3h}{16}\,\bk_1
            + \tfrac{9h}{16}\,\bk_4\right) \notag\\
  \bk_6 &= \bff\!\left(t_n + h,\;\by_n - \tfrac{3h}{7}\,\bk_1
            + \tfrac{2h}{7}\,\bk_2 + \tfrac{12h}{7}\,\bk_3
            - \tfrac{12h}{7}\,\bk_4 + \tfrac{8h}{7}\,\bk_5\right) \notag\\[6pt]
  \by_{n+1} &= \by_n + \frac{h}{90}\bigl(7\,\bk_1 + 32\,\bk_3
            + 12\,\bk_4 + 32\,\bk_5 + 7\,\bk_6\bigr)
  \label{eq:rk5}
\end{align}

\subsection{Butcher Tableau}
\[
\renewcommand{\arraystretch}{1.3}
\begin{array}{c|cccccc}
  0    &       &       &       &        &      &      \\
  1/4  & 1/4   &       &       &        &      &      \\
  1/4  & 1/8   & 1/8   &       &        &      &      \\
  1/2  & 0     & -1/2  & 1     &        &      &      \\
  3/4  & 3/16  & 0     & 0     & 9/16   &      &      \\
  1    & -3/7  & 2/7   & 12/7  & -12/7  & 8/7  &      \\
  \hline
       & 7/90  & 0     & 32/90 & 12/90  & 32/90 & 7/90 \\
\end{array}
\]

\subsection{Properties}
\begin{center}
\begin{tabular}{ll}
  \hline
  Order          & 5 \\
  Stages         & 6 \\
  Implicit       & No \\
  Symplectic     & No \\
  Local error    & $\bigO(h^6)$ \\
  \hline
\end{tabular}
\end{center}

%=============================================================================
\section{Bogacki--Shampine (RK23) --- \texttt{rk23\_step}}
\label{sec:rk23}
%=============================================================================

The Bogacki--Shampine method is a three-stage embedded Runge--Kutta pair of
orders~2 and~3. The implementation uses the third-order solution to advance
the state. The embedded pair structure makes it well suited for adaptive
step-size control (only the fixed-step advance is implemented).

\subsection{Formulation}
\begin{align}
  \bk_1 &= \bff(t_n,\;\by_n) \notag\\
  \bk_2 &= \bff\!\left(t_n + \tfrac{h}{2},\;\by_n + \tfrac{h}{2}\,\bk_1\right) \notag\\
  \bk_3 &= \bff\!\left(t_n + \tfrac{3h}{4},\;\by_n + \tfrac{3h}{4}\,\bk_2\right) \notag\\[6pt]
  \by_{n+1} &= \by_n + h\left(\tfrac{2}{9}\,\bk_1 + \tfrac{1}{3}\,\bk_2
               + \tfrac{4}{9}\,\bk_3\right)
  \label{eq:rk23}
\end{align}

\subsection{Butcher Tableau}
\[
\renewcommand{\arraystretch}{1.3}
\begin{array}{c|ccc}
  0    &      &      &     \\
  1/2  & 1/2  &      &     \\
  3/4  & 0    & 3/4  &     \\
  \hline
       & 2/9  & 1/3  & 4/9 \\
\end{array}
\]

\subsection{Properties}
\begin{center}
\begin{tabular}{ll}
  \hline
  Order (advance)  & 3 \\
  Order (embedded) & 2 \\
  Stages           & 3 \\
  Implicit         & No \\
  Symplectic       & No \\
  Local error      & $\bigO(h^4)$ \\
  \hline
\end{tabular}
\end{center}

%=============================================================================
\section{Runge--Kutta--Fehlberg (RKF45) --- \texttt{rkf45\_step}}
\label{sec:rkf45}
%=============================================================================

The Runge--Kutta--Fehlberg method is a six-stage embedded pair of orders~4
and~5. The fifth-order solution is used to advance the state, making it an
efficient high-accuracy integrator. Originally designed for adaptive step-size
control; the implementation uses fixed steps.

\subsection{Formulation}
\begin{align}
  \bk_1 &= \bff(t_n,\;\by_n) \notag\\
  \bk_2 &= \bff\!\left(t_n + \tfrac{h}{4},\;\by_n
            + \tfrac{h}{4}\,\bk_1\right) \notag\\
  \bk_3 &= \bff\!\left(t_n + \tfrac{3h}{8},\;\by_n
            + \tfrac{3h}{32}\,\bk_1 + \tfrac{9h}{32}\,\bk_2\right) \notag\\
  \bk_4 &= \bff\!\left(t_n + \tfrac{12h}{13},\;\by_n
            + \tfrac{1932h}{2197}\,\bk_1
            - \tfrac{7200h}{2197}\,\bk_2
            + \tfrac{7296h}{2197}\,\bk_3\right) \notag\\
  \bk_5 &= \bff\!\left(t_n + h,\;\by_n
            + \tfrac{439h}{216}\,\bk_1 - 8h\,\bk_2
            + \tfrac{3680h}{513}\,\bk_3
            - \tfrac{845h}{4104}\,\bk_4\right) \notag\\
  \bk_6 &= \bff\!\left(t_n + \tfrac{h}{2},\;\by_n
            - \tfrac{8h}{27}\,\bk_1 + 2h\,\bk_2
            - \tfrac{3544h}{2565}\,\bk_3
            + \tfrac{1859h}{4104}\,\bk_4
            - \tfrac{11h}{40}\,\bk_5\right) \notag\\[6pt]
  \by_{n+1} &= \by_n + h\left(\tfrac{16}{135}\,\bk_1
               + \tfrac{6656}{12825}\,\bk_3
               + \tfrac{28561}{56430}\,\bk_4
               - \tfrac{9}{50}\,\bk_5
               + \tfrac{2}{55}\,\bk_6\right)
  \label{eq:rkf45}
\end{align}

\subsection{Butcher Tableau}
\[
\renewcommand{\arraystretch}{1.3}
\begin{array}{c|cccccc}
  0      &            &          &              &            &        &      \\
  1/4    & 1/4        &          &              &            &        &      \\
  3/8    & 3/32       & 9/32     &              &            &        &      \\
  12/13  & 1932/2197  & -7200/2197 & 7296/2197  &            &        &      \\
  1      & 439/216    & -8       & 3680/513     & -845/4104  &        &      \\
  1/2    & -8/27      & 2        & -3544/2565   & 1859/4104  & -11/40 &      \\
  \hline
  \text{5th} & 16/135 & 0        & 6656/12825   & 28561/56430 & -9/50 & 2/55 \\
\end{array}
\]

\subsection{Properties}
\begin{center}
\begin{tabular}{ll}
  \hline
  Order (advance)  & 5 \\
  Order (embedded) & 4 \\
  Stages           & 6 \\
  Implicit         & No \\
  Symplectic       & No \\
  Local error      & $\bigO(h^6)$ \\
  \hline
\end{tabular}
\end{center}

%=============================================================================
\section{Semi-Implicit Euler (Symplectic 1st Order) --- \texttt{semi\_implicit\_euler\_step}}
\label{sec:euler}
%=============================================================================

The semi-implicit (symplectic) Euler method updates velocity first using the
current-state derivative, then advances position using the \emph{new} velocity.
This asymmetric coupling makes the method symplectic, preserving the phase-space
structure of Hamiltonian systems and providing long-term energy stability that
the standard explicit Euler lacks.

\subsection{Formulation}
\begin{align}
  \omega_{n+1} &= \omega_n + h\,f_\omega(t_n,\;\theta_n,\;\omega_n) \notag\\
  \theta_{n+1} &= \theta_n + h\,\omega_{n+1}
  \label{eq:sie}
\end{align}
where $f_\omega$ is the angular-acceleration component of the derivative.

\subsection{Properties}
\begin{center}
\begin{tabular}{ll}
  \hline
  Order          & 1 \\
  Stages         & 1 \\
  Implicit       & Semi-implicit \\
  Symplectic     & Yes \\
  Local error    & $\bigO(h^2)$ \\
  \hline
\end{tabular}
\end{center}

%=============================================================================
\section{Leapfrog / St\"{o}rmer--Verlet (Symplectic 2nd Order) --- \texttt{leapfrog\_step}}
\label{sec:leapfrog}
%=============================================================================

The leapfrog method (drift-kick-drift form of St\"{o}rmer--Verlet) is a
second-order symplectic integrator. It advances position by a half-step,
evaluates the force at the midpoint, updates velocity over the full step,
then completes the position half-step. It is time-reversible and preserves
phase-space volume.

\subsection{Formulation}
\begin{align}
  \theta_{n+1/2} &= \theta_n + \tfrac{h}{2}\,\omega_n \notag\\
  \omega_{n+1}   &= \omega_n + h\,f_\omega\!\left(t_n + \tfrac{h}{2},\;\theta_{n+1/2}\right) \notag\\
  \theta_{n+1}   &= \theta_{n+1/2} + \tfrac{h}{2}\,\omega_{n+1}
  \label{eq:leapfrog}
\end{align}

\subsection{Properties}
\begin{center}
\begin{tabular}{ll}
  \hline
  Order             & 2 \\
  Force evaluations & 1 per step \\
  Implicit          & No (explicit for separable~$H$) \\
  Symplectic        & Yes \\
  Time-reversible   & Yes \\
  Local error       & $\bigO(h^3)$ \\
  \hline
\end{tabular}
\end{center}

%=============================================================================
\section{Ruth's 4th-Order Symplectic Integrator --- \texttt{ruth4\_step}}
\label{sec:ruth4}
%=============================================================================

A fourth-order symplectic integrator based on the Yoshida/Forest--Ruth
decomposition. It splits the Hamiltonian flow into alternating position kicks
and momentum kicks with specially chosen coefficients that cancel lower-order
error terms while preserving the symplectic structure exactly.

\subsection{Coefficients}

Define $w = 2^{1/3}$ and $\alpha = \dfrac{1}{2 - w}$. The position
coefficients~$c_i$ and momentum coefficients~$d_i$ are:
\begin{align}
  c_1 &= \frac{\alpha}{2},        & d_1 &= \alpha, \notag\\
  c_2 &= \frac{(1-w)\,\alpha}{2}, & d_2 &= -w\,\alpha, \notag\\
  c_3 &= c_2,                     & d_3 &= \alpha, \notag\\
  c_4 &= c_1.                     &     &
  \label{eq:ruth-coeff}
\end{align}

These satisfy the consistency conditions $\sum_i c_i = 1$ and $\sum_i d_i = 1$.

\subsection{Formulation}

The integrator applies four position drifts interleaved with three momentum
kicks:
\begin{align}
  \theta &\leftarrow \theta + c_1\,h\,\omega \notag\\
  \omega &\leftarrow \omega + d_1\,h\,f_\omega(\theta) \notag\\
  \theta &\leftarrow \theta + c_2\,h\,\omega \notag\\
  \omega &\leftarrow \omega + d_2\,h\,f_\omega(\theta) \notag\\
  \theta &\leftarrow \theta + c_3\,h\,\omega \notag\\
  \omega &\leftarrow \omega + d_3\,h\,f_\omega(\theta) \notag\\
  \theta &\leftarrow \theta + c_4\,h\,\omega
  \label{eq:ruth4}
\end{align}

\subsection{Properties}
\begin{center}
\begin{tabular}{ll}
  \hline
  Order             & 4 \\
  Force evaluations & 3 (with 4 drift sub-steps) \\
  Implicit          & No (for separable Hamiltonians) \\
  Symplectic        & Yes \\
  Time-reversible   & Yes \\
  Local error       & $\bigO(h^5)$ \\
  \hline
\end{tabular}
\end{center}

%=============================================================================
\section{Velocity Verlet --- \texttt{velocity\_verlet\_step}}
\label{sec:vv}
%=============================================================================

The Velocity Verlet algorithm is a second-order method designed for the direct
second-order ODE $\ddot\theta = a(t,\theta)$. It is algebraically equivalent to
the St\"{o}rmer--Verlet method but is formulated to give synchronous position and
velocity output at each step. For conservative (velocity-independent) forces it
is symplectic.

\subsection{Formulation}
\begin{align}
  a_n         &= a(t_n,\;\theta_n) \notag\\
  \theta_{n+1} &= \theta_n + h\,\omega_n + \tfrac{h^2}{2}\,a_n \notag\\
  a_{n+1}     &= a(t_{n+1},\;\theta_{n+1}) \notag\\
  \omega_{n+1} &= \omega_n + \tfrac{h}{2}\bigl(a_n + a_{n+1}\bigr)
  \label{eq:vv}
\end{align}

\subsection{Properties}
\begin{center}
\begin{tabular}{ll}
  \hline
  Order             & 2 \\
  Force evaluations & 2 per step (1 reusable from previous step) \\
  Implicit          & No \\
  Symplectic        & Yes (conservative forces) \\
  Local error       & $\bigO(h^3)$ \\
  \hline
\end{tabular}
\end{center}

%=============================================================================
\section{Classical Runge--Kutta--Nystr\"{o}m (4th Order) --- \texttt{runge\_kutta\_nystrom\_step}}
\label{sec:rkn}
%=============================================================================

The Runge--Kutta--Nystr\"{o}m (RKN) method is a fourth-order integrator
specialised for second-order ODEs of the form $\ddot\theta = a(t,\theta)$. By
exploiting the second-order structure directly it avoids reducing the problem to
a first-order system, and produces both $\theta$ and~$\omega$ to fourth-order
accuracy.

\subsection{Formulation}
\begin{align}
  k_{1,\theta} &= \omega_n, &
  k_{1,\omega} &= a(t_n,\;\theta_n) \notag\\[4pt]
  k_{2,\theta} &= \omega_n + \tfrac{h}{2}\,k_{1,\omega}, &
  k_{2,\omega} &= a\!\left(t_n + \tfrac{h}{2},\;\theta_n + \tfrac{h}{2}\,k_{1,\theta}\right) \notag\\[4pt]
  k_{3,\theta} &= \omega_n + \tfrac{h}{2}\,k_{2,\omega}, &
  k_{3,\omega} &= a\!\left(t_n + \tfrac{h}{2},\;\theta_n + \tfrac{h}{2}\,k_{2,\theta}\right) \notag\\[4pt]
  k_{4,\theta} &= \omega_n + h\,k_{3,\omega}, &
  k_{4,\omega} &= a\!\left(t_n + h,\;\theta_n + h\,k_{3,\theta}\right) \notag\\[6pt]
  \theta_{n+1} &= \theta_n + \frac{h}{6}\bigl(k_{1,\theta} + 2\,k_{2,\theta}
                   + 2\,k_{3,\theta} + k_{4,\theta}\bigr) \notag\\[4pt]
  \omega_{n+1} &= \omega_n + \frac{h}{6}\bigl(k_{1,\omega} + 2\,k_{2,\omega}
                   + 2\,k_{3,\omega} + k_{4,\omega}\bigr)
  \label{eq:rkn}
\end{align}

\subsection{Properties}
\begin{center}
\begin{tabular}{ll}
  \hline
  Order             & 4 \\
  Force evaluations & 4 \\
  Implicit          & No \\
  Local error       & $\bigO(h^5)$ \\
  \hline
\end{tabular}
\end{center}

%=============================================================================
\section{Numerov Method (Implicit, 4th Order) --- \texttt{numerov\_next\_theta}}
\label{sec:numerov}
%=============================================================================

Numerov's method is a high-accuracy multistep scheme for second-order ODEs with
no first-derivative term, $\ddot\theta = a(t,\theta)$. It achieves fourth-order
accuracy using only three consecutive position values and their accelerations.
The method is implicit because $a_{n+1}$ depends on the unknown~$\theta_{n+1}$;
the implementation resolves this with Newton iteration using a numerically
estimated Jacobian.

\subsection{Recurrence}
\begin{equation}\label{eq:numerov}
  \theta_{n+1} - 2\,\theta_n + \theta_{n-1}
  = \frac{h^2}{12}\bigl(a_{n-1} + 10\,a_n + a_{n+1}\bigr),
\end{equation}
where $a_k = a(t_k,\theta_k)$.

\subsection{Newton Iteration}

Since $a_{n+1} = a(t_{n+1},\theta_{n+1})$ depends on the unknown, the
recurrence is rearranged into the root-finding problem
\begin{equation}\label{eq:numerov-root}
  F(\theta_{n+1}) = \theta_{n+1}
  - \underbrace{\left[2\,\theta_n - \theta_{n-1}
    + \frac{h^2}{12}\bigl(a_{n-1} + 10\,a_n\bigr)\right]}_{\text{RHS}}
  - \frac{h^2}{12}\,a(t_{n+1},\theta_{n+1}) = 0.
\end{equation}

The Jacobian is approximated via central finite differences:
\begin{equation}\label{eq:numerov-jac}
  \frac{\partial a}{\partial\theta}\bigg|_{\theta_{n+1}}
  \approx
  \frac{a(t_{n+1},\,\theta_{n+1}+\varepsilon) - a(t_{n+1},\,\theta_{n+1}-\varepsilon)}{2\varepsilon},
\end{equation}
and the Newton update is
\begin{equation}\label{eq:numerov-update}
  \theta_{n+1} \leftarrow \theta_{n+1}
  - \frac{F(\theta_{n+1})}{1 - \tfrac{h^2}{12}\,\partial a / \partial\theta}.
\end{equation}

\subsection{Bootstrap and Velocity Recovery}

Numerov requires two prior position values. The implementation bootstraps
$\theta_1$ and $\theta_2$ using the 4th-order Runge--Kutta--Nystr\"{o}m method
(Section~\ref{sec:rkn}). All arithmetic during the Numerov phase is performed in
\texttt{long double} for enhanced precision.

Velocities are recovered \emph{a~posteriori} from the position trajectory using
the corrected central-difference formula:
\begin{equation}\label{eq:numerov-vel}
  \omega_i = \frac{\theta_{i+1} - \theta_{i-1}}{2h}
             - \frac{h}{12}\bigl(a_{i+1} - a_{i-1}\bigr),
  \qquad 1 \le i \le N-1,
\end{equation}
and the endpoint velocity uses a five-point backward-difference stencil:
\begin{equation}\label{eq:numerov-endpoint}
  \omega_N = \frac{25\,\theta_N - 48\,\theta_{N-1} + 36\,\theta_{N-2}
             - 16\,\theta_{N-3} + 3\,\theta_{N-4}}{12h}.
\end{equation}

\subsection{Properties}
\begin{center}
\begin{tabular}{ll}
  \hline
  Order     & 4 (global) \\
  Type      & Multistep (requires two prior positions) \\
  Implicit  & Yes (Newton iteration per step) \\
  Precision & \texttt{long double} arithmetic \\
  Local error & $\bigO(h^6)$ \\
  \hline
\end{tabular}
\end{center}

%=============================================================================
\section{Discrete Exponential Nystr\"{o}m (DEN3) --- \texttt{den3\_step}}
\label{sec:den3}
%=============================================================================

The DEN3 method is a third-order exponential integrator designed for the damped
nonlinear oscillator equation
\begin{equation}\label{eq:den-ode}
  \ddot\theta + 2\gamma\,\dot\theta + \omega_0^2\,\theta = g(t,\theta,\dot\theta),
\end{equation}
where $\gamma$ is the damping coefficient, $\omega_0$ is the natural frequency,
and $g$ is the nonlinear residual. The key idea is to solve the linear part
$\ddot\theta + 2\gamma\dot\theta + \omega_0^2\theta = 0$ \emph{exactly} via a
matrix exponential propagator, and treat the nonlinear residual with a
three-stage quadrature of Simpson type.

\subsection{Propagator Matrix}

The state vector $\bx = (\theta,\;\omega)^T$ satisfies the homogeneous equation
$\dot\bx = A\bx$ with
\begin{equation}\label{eq:den-matrix}
  A = \begin{pmatrix} 0 & 1 \\ -\omega_0^2 & -2\gamma \end{pmatrix}.
\end{equation}

The matrix exponential $\Phi(\tau) = e^{A\tau}$ is computed analytically in
three regimes depending on the discriminant
$\Delta = \omega_0^2 - \gamma^2$:

\paragraph{Underdamped ($\Delta > 0$).}
Let $\omega_d = \sqrt{\omega_0^2 - \gamma^2}$.
\begin{equation}\label{eq:phi-under}
  \Phi(\tau) = e^{-\gamma\tau}
  \begin{pmatrix}
    \cos\omega_d\tau + \dfrac{\gamma}{\omega_d}\sin\omega_d\tau &
    \dfrac{1}{\omega_d}\sin\omega_d\tau \\[8pt]
    -\dfrac{\omega_0^2}{\omega_d}\sin\omega_d\tau &
    \cos\omega_d\tau - \dfrac{\gamma}{\omega_d}\sin\omega_d\tau
  \end{pmatrix}
\end{equation}

\paragraph{Overdamped ($\Delta < 0$).}
Let $\omega_h = \sqrt{\gamma^2 - \omega_0^2}$.
\begin{equation}\label{eq:phi-over}
  \Phi(\tau) = e^{-\gamma\tau}
  \begin{pmatrix}
    \cosh\omega_h\tau + \dfrac{\gamma}{\omega_h}\sinh\omega_h\tau &
    \dfrac{1}{\omega_h}\sinh\omega_h\tau \\[8pt]
    -\dfrac{\omega_0^2}{\omega_h}\sinh\omega_h\tau &
    \cosh\omega_h\tau - \dfrac{\gamma}{\omega_h}\sinh\omega_h\tau
  \end{pmatrix}
\end{equation}

\paragraph{Critically damped ($\Delta = 0$).}
\begin{equation}\label{eq:phi-crit}
  \Phi(\tau) = e^{-\gamma\tau}
  \begin{pmatrix}
    1 + \gamma\tau & \tau \\
    -\omega_0^2\,\tau & 1 - \gamma\tau
  \end{pmatrix}
\end{equation}

\subsection{Constant-Residual Response}

The particular solution for a constant unit forcing $g = 1$ from zero initial
conditions is:
\begin{equation}\label{eq:den-response}
  \br(\tau) =
  \begin{pmatrix}
    r_\theta(\tau) \\ r_\omega(\tau)
  \end{pmatrix}
  =
  \begin{pmatrix}
    \dfrac{1 - \Phi_{22}(\tau) - 2\gamma\,\Phi_{12}(\tau)}{\omega_0^2} \\[8pt]
    \Phi_{12}(\tau)
  \end{pmatrix},
\end{equation}
which follows from the identity
$\int_0^\tau \Phi_{12}(s)\,ds = r_\theta(\tau)$ and the relation
$r_\omega(\tau) = \dot{r}_\theta(\tau) = \Phi_{12}(\tau)$.

\subsection{Three-Stage Stepping Procedure}

Given state $(\theta_n,\omega_n)$ at time~$t_n$:

\begin{enumerate}
  \item \textbf{Stage~1}: evaluate residual at the current state
  \[
    g_1 = g(t_n,\;\theta_n,\;\omega_n).
  \]

  \item \textbf{Stage~2}: propagate a half-step with constant forcing~$g_1$
  \[
    \begin{pmatrix}\theta_2 \\ \omega_2\end{pmatrix}
    = \Phi\!\left(\tfrac{h}{2}\right)
      \begin{pmatrix}\theta_n \\ \omega_n\end{pmatrix}
    + g_1\,\br\!\left(\tfrac{h}{2}\right),
  \]
  then evaluate $g_2 = g(t_n + h/2,\;\theta_2,\;\omega_2)$.

  \item \textbf{Stage~3}: propagate a full step with constant forcing~$g_2$
  \[
    \begin{pmatrix}\theta_3 \\ \omega_3\end{pmatrix}
    = \Phi(h)\begin{pmatrix}\theta_n \\ \omega_n\end{pmatrix}
    + g_2\,\br(h),
  \]
  then evaluate $g_3 = g(t_n + h,\;\theta_3,\;\omega_3)$.

  \item \textbf{Final combination} using Simpson-type weights
        $\bigl(\tfrac{1}{6},\;\tfrac{2}{3},\;\tfrac{1}{6}\bigr)$:
  \begin{align}
    \theta_{n+1} &= \Phi_{11}\,\theta_n + \Phi_{12}\,\omega_n
      + h\left[\Phi_{12}(h)\,\frac{g_1}{6}
        + \Phi_{12}\!\left(\tfrac{h}{2}\right)\frac{2g_2}{3}\right]
      \label{eq:den3-theta}\\[4pt]
    \omega_{n+1} &= \Phi_{21}\,\theta_n + \Phi_{22}\,\omega_n
      + h\left[\Phi_{22}(h)\,\frac{g_1}{6}
        + \Phi_{22}\!\left(\tfrac{h}{2}\right)\frac{2g_2}{3}
        + \frac{g_3}{6}\right]
      \label{eq:den3-omega}
  \end{align}
\end{enumerate}

\subsection{Properties}
\begin{center}
\begin{tabular}{ll}
  \hline
  Order             & 3 \\
  Stages            & 3 residual evaluations \\
  Implicit          & No \\
  Exponential       & Yes (exact linear propagation) \\
  Suited for        & Damped oscillators with nonlinear forcing \\
  Local error       & $\bigO(h^4)$ \\
  \hline
\end{tabular}
\end{center}

%=============================================================================
\section{Summary}
\label{sec:summary}
%=============================================================================

\begin{center}
\small
\renewcommand{\arraystretch}{1.2}
\begin{tabular}{>{\raggedright}p{4.2cm} l c c c >{\raggedright\arraybackslash}p{2.5cm}}
  \hline
  \textbf{Method} & \textbf{Function} & \textbf{Order} & \textbf{Symplectic} & \textbf{Implicit} & \textbf{Family} \\
  \hline
  RK3                      & \texttt{rk3\_step}                 & 3 & No  & No   & Runge--Kutta \\
  RK4                      & \texttt{rk4\_step}                 & 4 & No  & No   & Runge--Kutta \\
  RK5 (Butcher)            & \texttt{rk5\_step}                 & 5 & No  & No   & Runge--Kutta \\
  Bogacki--Shampine (RK23) & \texttt{rk23\_step}                & 3 & No  & No   & Embedded RK \\
  RKF45                    & \texttt{rkf45\_step}               & 5 & No  & No   & Embedded RK \\
  Semi-Implicit Euler      & \texttt{semi\_implicit\_euler\_step} & 1 & Yes & Semi & Symplectic \\
  Leapfrog                 & \texttt{leapfrog\_step}            & 2 & Yes & No   & Symplectic \\
  Ruth4                    & \texttt{ruth4\_step}               & 4 & Yes & No   & Symplectic \\
  Velocity Verlet          & \texttt{velocity\_verlet\_step}    & 2 & Yes & No   & Verlet \\
  Runge--Kutta--Nystr\"{o}m & \texttt{runge\_kutta\_nystrom\_step} & 4 & No & No & Nystr\"{o}m \\
  Numerov                  & \texttt{numerov\_next\_theta}      & 4 & No  & Yes  & Multistep \\
  DEN3                     & \texttt{den3\_step}                & 3 & No  & No   & Exponential \\
  \hline
\end{tabular}
\end{center}

\end{document}
